% !TeX root = ../main.tex
\chapter{Design Calculations and Circuit Analysis}
\label{app:design-calculations}
This appendix contains the electrical calculations used to select and verify the
preliminary hardware. In addition to nominal values, voltage, current, power,
logic-level compatibility, and component capability are considered so that the
design does not intentionally operate a component beyond its specified limits.

\section{Current-Sense Shunt Resistor Sizing}
\label{sec:shunt-sizing}
The LTC2945 measures the differential voltage between SENSE+ and SENSE-- with a
full-scale range of \SI{102.4}{\milli\volt}. The ideal shunt resistance for a
specified maximum measurable current is therefore
\begin{equation}
    R_{SHUNT,ideal}=\frac{V_{SENSE,FS}}{I_{MAX}}
    =\frac{\SI{102.4}{\milli\volt}}{I_{MAX}}.
\end{equation}
To prevent saturation, a scaling factor, $K_H$, is applied to leave some headroom:
\begin{equation}
    R_{SHUNT}=K_H\frac{\SI{102.4}{\milli\volt}}{I_{MAX}},
    \qquad 0<K_H<1.
\end{equation}
For example, $K_H=0.80$ reserves 20\% headroom for overload, tolerance, and
transient conditions.

The shunt power at maximum continuous current is
\begin{equation}
    P_{SHUNT}=I_{MAX}^{2}R_{SHUNT}.
\end{equation}
The installed resistor power rating should exceed this value.

\subsection{Current Measurement Resolution}
The LTC2945 current-sense ADC has a nominal sense-voltage resolution of
\SI{25}{\micro\volt} per LSB. The current represented by one ADC count is
\begin{equation}
    I_{LSB}=\frac{\SI{25}{\micro\volt}}{R_{SHUNT}}.
\end{equation}
This relationship provides an important design tradeoff: a smaller shunt lowers
power dissipation but also increases the minimum current increment represented by
one ADC count.

\subsection{Production Shunt Design Table}
\begin{table}[H]
    \centering
    \caption{Production current-sense resistor calculations.}
    \label{tab:production-shunts}
    \begin{tabular}{lrrrrrr}
        \toprule
        Rail & $I_{MAX}$ & $R_{ideal}$ & Selected $R$ &
        $V_{SHUNT}$ & $I_{LSB}$ & $P_{MAX}$ \\
        & (A) & ($\Omega$) & ($\Omega$) &
        (mV) & (mA/LSB) & (W) \\
        \midrule
        +5VA  & 2.00 & 0.0512 & 0.039 & 78.0 & 0.641 & 0.156 \\
        +5VB  & 2.00 & 0.0512 & 0.039 & 78.0 & 0.641 & 0.156 \\
        +5V   & 1.00 & 0.1024 & 0.050 & 50.0 & 0.500 & 0.050 \\
        -5V   & 1.00 & 0.1024 & 0.050 & 50.0 & 0.500 & 0.050 \\
        +15V  & 1.00 & 0.1024 & 0.050 & 50.0 & 0.500 & 0.050 \\
        -15V  & 1.00 & 0.1024 & 0.050 & 50.0 & 0.500 & 0.050 \\
        \bottomrule
    \end{tabular}
\end{table}

\subsection{Evaluation-Board Shunt Calculation}
The prototype LTC2945 evaluation board was tested from a \SI{5}{\volt} source
with a load range of approximately \SI{1}{\kilo\ohm} to \SI{10}{\kilo\ohm}.
This is a low-current bench-test condition and is separate from the production
shunt design.

At the minimum load resistance, the shunt required to place the current-sense
input at full scale is
\begin{equation}
    R_{S,FS}=\frac{V_{SHUNT}R_L}{V_S-V_{SHUNT}}
    =\frac{(\SI{0.1024}{\volt})(\SI{1}{\kilo\ohm})}
    {\SI{5}{\volt}-\SI{0.1024}{\volt}}
    =\SI{20.91}{\ohm}.
\end{equation}
Applying approximately 15\% headroom gives
\begin{equation}
    R_S=(0.85)(\SI{20.91}{\ohm})\approx\SI{17.8}{\ohm}.
\end{equation}
At a \SI{10}{\kilo\ohm} load,
\begin{equation}
    I_{MIN}=\frac{\SI{5}{\volt}}
    {\SI{10}{\kilo\ohm}+\SI{17.8}{\ohm}}
    \approx\SI{499.1}{\micro\ampere},
\end{equation}
and therefore
\begin{equation}
    V_{SHUNT,MIN}=I_{MIN}R_S
    \approx(\SI{499.1}{\micro\ampere})(\SI{17.8}{\ohm})
    \approx\SI{8.88}{\milli\volt}.
\end{equation}
The prototype sense voltage therefore remains inside the LTC2945
\SI{102.4}{\milli\volt} full-scale range over the intended test loads.

\section{PCB Trace and Copper-Pour Current Capacity}
\label{sec:trace-width}
Power-path copper must carry the maximum rail current without excessive
temperature rise or voltage drop. KiCad's track-width calculator uses a
closed-form fit to IPC-2152 data. For an external trace it can be expressed as
\begin{equation}
    \Delta T = K I^a W^b T_h^c,
\end{equation}
where $I$ is current in amperes, $W$ is trace width in mils, $T_h$ is copper
thickness in mils, and for external tracks the KiCad coefficients are
\begin{equation}
    K=215.3,\qquad a=2,\qquad b=-1.15,\qquad c=-1.0.
\end{equation}
Solving for minimum width,
\begin{equation}
    W_{MIN}=\left(\frac{\Delta T}
    {K I^a T_h^c}\right)^{1/b}.
\end{equation}
For \SI{1}{oz} copper ($T_h\approx\SI{1.378}{mil}$) and a preliminary
\SI{10}{\celsius} allowed rise, the calculated external-trace widths are
approximately \SI{3.27}{mil} at \SI{0.5}{A}, \SI{10.9}{mil} at \SI{1}{A}, and
\SI{36.4}{mil} at \SI{2}{A}. These are thermal minimums only; the actual layout
should use wider traces or copper pours where space permits.

\section{I\texorpdfstring{\textsuperscript{2}}{2}C Pull-Up Resistor Sizing}
\label{sec:i2c-pullups}
The I2C bus buffer and I2C line driver both have internal active current sources that eliminate the need for external pullups. Footprints were placed for 4.7k pullups but are set as Do Not Populate.

\section{LTC2945 Voltage Measurement Resolution}
The LTC2945 rail-voltage channel has a nominal full-scale range of
\SI{102.4}{\volt} and a nominal resolution of \SI{25}{\milli\volt} per LSB.
The quantization-only percentage of a measured rail voltage is therefore
\begin{equation}
    \epsilon_Q=\frac{\SI{25}{\milli\volt}}{V_{RAIL}}\times100\%.
\end{equation}
For example, one count corresponds to approximately 0.50\% of a \SI{5}{\volt}
rail and 0.17\% of a \SI{15}{\volt} rail.
This is only ADC resolution and not total measurement
accuracy.

\section{Negative-Rail Sign Convention}
The -5 V and -15 V supplies are electrically isolated from the controller. The
monitor can therefore be placed in the high side of each isolated rail and read a
positive local differential quantity even when the system-level rail is labeled
negative. Firmware applies the configured sign after acquisition:
\begin{equation}
    V_{DISPLAY}=S_VV_{MEASURED},\qquad S_V\in\{+1,-1\}.
\end{equation}
The same channel metadata is used by the local display and network interface so
that a negative rail is represented consistently throughout the system.

\section{Alarm-Limit Scaling}
User alarm thresholds are stored and compared in engineering units. For a lower
voltage limit $V_L$ and upper voltage limit $V_H$, the normal voltage state is
\begin{equation}
    V_L\le V_{MEAS}\le V_H.
\end{equation}
Similarly, for lower and upper current limits,
\begin{equation}
    I_L\le I_{MEAS}\le I_H.
\end{equation}
If raw ADC thresholds are ever written directly to the LTC2945, the corresponding
register value must be calculated from the ADC LSB size. For the current-sense
channel,
\begin{equation}
    N_I=\frac{I_{LIMIT}R_{SHUNT}}{\SI{25}{\micro\volt/LSB}},
\end{equation}
and for the rail-voltage channel,
\begin{equation}
    N_V=\frac{V_{LIMIT}}{\SI{25}{\milli\volt/LSB}}.
\end{equation}

\section{Controller-Supply Power Budget}

A preliminary controller power budget was developed using datasheet maximum
values where available and conservative design-current allowances for assemblies
where a complete board-level maximum current is not specified. A 25\% design
margin is applied to the calculated load.

\begin{table}[H]
    \centering
    \caption{Preliminary controller-side power budget.}
    \label{tab:power-budget}
    \begin{tabular}{lrrr}
        \toprule
        Load & Supply & Design Current & Power \\
        & (V) & (A) & (W) \\
        \midrule
        Arduino GIGA R1 WiFi
        & 5.0 & 0.500 & 2.500 \\

        GIGA Display Shield
        & 3.3 & 0.250 & 0.825 \\

        W6100 Ethernet interface
        & 3.3 & 0.265 & 0.875 \\

        I\textsuperscript{2}C isolation/controller circuitry
        & 3.3 & 0.065 & 0.215 \\
        \midrule
        Subtotal
        & -- & -- & 4.415 \\

        25\% design margin
        & -- & -- & 1.104 \\
        \midrule
        \textbf{Total}
        & -- & -- & \textbf{5.519} \\
        \bottomrule
    \end{tabular}
\end{table}



\section{Component Capability Checks}
\label{sec:component-capability}
The calculations above establish the required electrical values. The following
checks compare those values with the capability of the selected components.

\section{Current-Sense Resistor Power Margin}
The selected LVK-series 39 m\(\Omega\) and 50 m\(\Omega\) four-terminal shunts
are rated at 0.75 W. At the maximum design current of 2 A, the 39 m\(\Omega\)
resistor dissipates
\begin{equation}
    P=(2\,\mathrm{A})^2(0.039\,\Omega)=0.156\,\mathrm{W}.
\end{equation}
This is approximately 20.8\% of the 0.75 W rating, corresponding to about a
4.8:1 power-rating margin before considering PCB temperature and derating. At
1 A, a 50 m\(\Omega\) shunt dissipates only 0.050 W, approximately 6.7\% of a
0.75 W rating.

\section{LTC2945 Input and ADC Range}
The highest monitored nominal rail is 15 V, far below the LTC2945 0--80 V rail
monitoring range. The selected production shunts produce 78 mV at 2 A on the
39 m\(\Omega\) channels and 50 mV at 1 A on the 50 m\(\Omega\) channels. Both
values are below the 102.4 mV current-sense full scale, leaving approximately
24.4 mV and 52.4 mV of ADC headroom respectively. The LTC2945 can operate from
an external 2.7--5.9 V supply, so a 3.3 V sensor-side logic supply is within its
specified operating range.

\section{Microcontroller I/O Compatibility}
The Arduino GIGA R1 WiFi operates with 3.3 V logic and specifies 8 mA maximum DC
current per I/O pin. The controller-side I\textsuperscript{2}C interface is a
logic interface, not a load-power interface. The TCA9803 used in the design
supports a 0.8--3.6 V A-side supply and a 1.65--3.6 V B-side supply and provides
an internal B-side current source of up to 3 mA. A 3 mA bus-low current is below
the 8 mA GIGA per-pin limit, providing more than a factor of two current margin.
The GIGA pins therefore do not directly source the optocoupler LED or rail-load
current; the external interface circuitry performs the required buffering and
isolation.

\section{Digital-Isolation Margin}
The ACPL-064L digital optocoupler is specified for 3.3 V/5 V operation and is
UL 1577 recognized for 3750 V\textsubscript{RMS} isolation for one minute. The
largest normal potential difference between the nominal monitored rails is only
30 V from -15 V to +15 V. The optocoupler rating therefore greatly exceeds the
functional isolation voltage required by this low-voltage monitoring system.
This comparison verifies component capability; it is not intended to redefine
the safety rating of the complete assembled product, which also depends on PCB
creepage, clearance, connectors, wiring, and enclosure construction.

\section{3.3 V Regulator Loading}
The TPS62912 is rated for up to 2 A output. Using the Week 1 design-current
allowances for the display, Ethernet interface, and controller-side isolation
circuitry gives an estimated 3.3 V load of
\begin{equation}
    I_{3.3}=0.250+0.265+0.065=0.580\,\mathrm{A}.
\end{equation}
This is approximately 29\% of the regulator's 2 A output rating and leaves
substantial current margin for transient and implementation variation.

\section{Tolerance and Measurement-Uncertainty Considerations}
The production current-sense resistors have a 1\% resistance tolerance. Because
current is calculated from \(I=V_{SHUNT}/R_{SHUNT}\), an uncorrected 1\% shunt
error produces approximately a 1\% gain error in calculated current. On a 2 A
range, this corresponds to approximately 20 mA, which is much larger than the
0.641 mA quantization step of the 39 m\(\Omega\) channel. The LTC2945 also has
its own conversion error; therefore component tolerance rather than ADC
quantization is an important contributor to absolute current accuracy. If
higher accuracy is required, the actual shunt resistance can be measured and a
per-channel calibration factor stored in firmware.

A conservative worst-case linear addition of a 1\% shunt tolerance and the
LTC2945's less-than-\(\pm0.75\%\) total unadjusted conversion error gives an
upper-bound estimate of approximately \(\pm1.75\%\) before wiring, thermal, and
calibration effects. This estimate is used for design margin rather than as a
claim of final calibrated system accuracy.

\section{Reference Specifications Used for Capability Checks}
\begin{itemize}
    \item Analog Devices, LTC2945 data sheet and product information: \url{https://www.analog.com/en/products/ltc2945.html}
    \item Arduino, GIGA R1 WiFi technical specifications: \url{https://store.arduino.cc/products/giga-r1-wifi}
    \item Broadcom, ACPL-064L low-power digital optocoupler data sheet/product page: \url{https://www.broadcom.com/products/optocouplers/industrial-plastic/digital-optocouplers/10mbd/acpl-064l-000e}
    \item Texas Instruments, TCA9803 I\textsuperscript{2}C buffer data sheet: \url{https://www.ti.com/product/TCA9803}
    \item Texas Instruments, TPS62912 buck converter data sheet: \url{https://www.ti.com/product/TPS62912}
    \item Ohmite LVK-series current-sense resistor specifications:
    \url{https://www.ohmite.com/assets/files/Datasheets/lvk-series.pdf}
\end{itemize}
